\section{Multiply residues without a 256-bit division}
The full product routine established how carries work. Reducing its 256-bit output
would require a wider reducer than the 128-bit interface just built. A different
multiplication algorithm avoids storing that full product: distribute the multiplier's
bits over repeated modular additions.

Initialize \(R=0\), \(X=A\bmod N\), and a private copy \(Y=B\). While \(Y\ne0\),
add \(X\) to \(R\) modulo \(N\) if \(Y\)'s low bit is one. Shift \(Y\) right.
If more bits remain, double \(X\) modulo \(N\). The invariant is
\(R+XY\equiv AB\pmod N\), with \(0\leq R,X<N\). Thus every modular-addition call
receives valid residue operands, including when \(A\) or \(B\) originally exceeded
\(N\).

For \(13\cdot11\bmod17\), the low-to-high multiplier bits select these transitions.
The final iteration skips doubling \(X\).
\tagpdfsetup{table/header-rows={1}}
\begin{center}
  \begin{tabular}{rrrrrr}
    \(R\) before & \(X\) before & \(Y\) before & \(R\) after & \(X\) after & \(Y\) after \\
    \hline
    0            & 13           & 11           & 13          & 9           & 5           \\
    13           & 9            & 5            & 5           & 1           & 2           \\
    5            & 1            & 2            & 5           & 2           & 1           \\
    5            & 2            & 1            & 7           & 2           & 0           \\
  \end{tabular}
\end{center}

A shift of the whole \(Y\) must transfer bits between words. Shifting each limb
independently drops those transferred bits. With a private, eight-byte-aligned copy
at stack offsets sixteen through thirty-one, two doubleword loads provide another
way to express the transfer. The low bit of the high doubleword becomes bit 63 of
the low doubleword.

\Needspace{14\baselineskip}
\MARSCodeLabel{Shift the private 128-bit multiplier right by one.}
\begin{minted}[samepage=true]{mars-idea}
ld $t0, 16($sp)
    ld $t1, 24($sp)
    dsrl $t0, $t0, 1
    dsll $t2, $t1, 63
    or $t0, $t0, $t2
    dsrl $t1, $t1, 1
    sd $t0, 16($sp)
    sd $t1, 24($sp)
\end{minted}

The same shift can be written using four word loads and carries between adjacent
limbs. Both forms must include every input bit through bit 127. Keep the caller's
multiplier unchanged by shifting only the private copy. An all-zero multiplier returns
zero; a multiplier with only bit 127 set must still reach its selected addition.

The result buffer is disjoint from caller inputs here. Use private \(R\) and \(X\)
for nested in-place additions, then copy the final four words to the caller's output.
This separates the public ownership rule from the internal update strategy.
